% Generated draft: requires mathematical and source review.
\documentclass[11pt]{article}
\usepackage[margin=25mm]{geometry}
\usepackage{fontspec}
\setmainfont{Latin Modern Roman}
\usepackage{amsmath,amssymb,mathtools,unicode-math}
\setmathfont{Latin Modern Math}
\usepackage{enumitem,needspace,xurl,hyperref}
\setlength{\parindent}{0pt}
\setlength{\parskip}{5pt}
\setlength{\emergencystretch}{3em}
\newcommand{\N}{\mathbb{N}}
\newcommand{\Z}{\mathbb{Z}}
\newcommand{\Q}{\mathbb{Q}}
\newcommand{\R}{\mathbb{R}}
\newcommand{\C}{\mathbb{C}}
\newcommand{\F}{\mathbb{F}}
\newcommand{\A}{\mathbb{A}}
\newcommand{\PP}{\mathbb{P}}
\newcommand{\E}{\mathbb{E}}
\newcommand{\eps}{\varepsilon}
\newcommand{\dd}{\,\mathrm{d}}
\begin{document}
\section*{UnsolvedMath ID: 800681}
\textbf{Source problem code:} catalogue entry
\subsection{Definitions and mathematical statement}
A holomorphic foliation or lamination is called minimal when every leaf is dense in its underlying foliated or laminated space. The supplied record does not specify the leaf dimension, whether singularities are permitted, or the precise ambient and support conventions for a lamination.

A current \(T\) is \(dd^c\)-closed if \(dd^cT=0\), and is non-closed if \(dT\ne0\). Here \(dd^c\) may be taken with any fixed standard nonzero normalization, which does not affect its vanishing. A current directed by a foliation or lamination is understood to be tangent to its leaves, with the precise definition depending on the current degree and the treatment of singularities.

To interpret extremality, suppose that \(\mathcal C_{\mathcal H}\) denotes the intended convex cone of directed \(dd^c\)-closed currents associated with a foliation or lamination \(\mathcal H\). A nonzero \(T\in\mathcal C_{\mathcal H}\) spans an extreme ray if
\[
 T=A+B,\quad A,B\in\mathcal C_{\mathcal H}
 \quad\Longrightarrow\quad
 A=aT,\quad B=bT\quad\text{for some }a,b\ge0.
\]
The source does not specify this cone, positivity, the degree of its currents, or a normalization defining an alternative convex set of currents. Consequently, this cone interpretation is conditional.

\textbf{Requested construction, under the extreme-ray interpretation.} Construct a compact complex manifold \(X\) and a minimal holomorphic foliation \(\mathcal H\) on \(X\) for which there exist an uncountable index set \(I\) and currents \((T_i)_{i\in I}\) such that
\[
\begin{gathered}
0\ne T_i\in\mathcal C_{\mathcal H},\qquad
 dd^cT_i=0,\qquad dT_i\ne0
 \quad(i\in I),\\
\R_{\ge0}T_i\text{ is an extreme ray of }\mathcal C_{\mathcal H}
 \quad(i\in I),\\
\R_{\ge0}T_i\ne\R_{\ge0}T_j\quad(i\ne j).
\end{gathered}
\]
The source separately asks for the corresponding construction with a minimal holomorphic lamination in place of the foliation. Requiring distinct rays expresses an interpretation of uncountably many extremal currents that avoids counting scalar multiples of a single current. The supplied record does not explicitly select this interpretation or supply the missing current conventions; a fully determined formulation requires clarification of those points.
\subsection{Short English statement}
Construct a minimal holomorphic foliation, and respectively a minimal holomorphic lamination, on a compact complex manifold carrying uncountably many extremal directed currents that satisfy \(dd^cT=0\) but \(dT\ne0\). The supplied statement leaves the current class and meaning of extremality unspecified; the displayed formulation conditionally interprets extremality through distinct extreme rays.
\subsection{Sources}
\begin{itemize}
\item Local UnsolvedMath registry, permanent ID \(800681\), exact source problem code \texttt{catalogue entry}, title: \emph{Holomorphic Foliation Theory --- Problem 8}. The registry attributes the problem to Dinh and Sibony, with proposed year 2018.
\item Dinh and Sibony, \emph{Some Open Problems on Holomorphic Foliation Theory} (2018), Problem 8. Supplied source URL: \url{https://arxiv.org/abs/1807.05640}. The record reports access on 2026-10-16 and extraction by \texttt{arxiv\_latex}; these are supplied provenance statements, not an independently verified visit or extraction.
\item The supplied registry labels the problem \texttt{open} and records an update on 2026-10-16. This is a supplied status label, not an independently verified assessment of the literature.
\end{itemize}
\end{document}
